% titel: Blatt 4, A3: Wellenlänge, Frequenz, Wellenzahl in Li2+
% Teil II: der durchgerechnete Weg.

\subsection*{Gegeben und gesucht}
Übergang $n = 5 \rightarrow n = 4$ im Ion Li$^{2+}$ ($Z=3$). Gesucht:
Wellenlänge $\lambda$, Frequenz $\nu$, Wellenzahl $\tilde\nu$.

\subsection*{Ansatz}
Wasserstoffähnliche Formel $\nu = R_\infty\,Z^2(1/n_f^2 - 1/n_i^2)$; daraus
$\lambda = c/\nu$ und $\tilde\nu = 1/\lambda$.

\subsection*{Rechnung}
\begin{align*}
  \nu &= 3.29\times10^{15}\times 3^2
       \left(\frac{1}{4^2} - \frac{1}{5^2}\right)
  = 3.29\times10^{15}\times 9 \times 0.0225
  = 6.66\times10^{14}\ \mathrm{Hz} \\[2pt]
  \lambda &= \frac{2.998\times10^{8}}{6.66\times10^{14}}
  = 4.50\times10^{-7}\ \mathrm{m}
  = \qty{450}{nm} \\[2pt]
  \tilde\nu &= \frac{1}{4.50\times10^{-7}}
  = 2.22\times10^{6}\ \mathrm{m^{-1}}
  = \qty{22200}{cm^{-1}}.
\end{align*}

\begin{ergebnis}
$\lambda \approx \qty{450}{nm}$,\quad
$\nu \approx \qty{6.66e14}{s^{-1}}$,\quad
$\tilde\nu \approx \qty{2.22e6}{m^{-1}} = \qty{22200}{cm^{-1}}$.
\end{ergebnis}

\begin{kontrolle}
Der gleiche Übergang liegt beim H ($Z=1$) bei $\lambda \approx 4050$ nm
(Infrarot); mit $Z=3$ skalieren Frequenz und Wellenzahl um $Z^2=9$, die
Wellenlänge um $1/9$: $4050/9 \approx 450$ nm --- konsistent. Die Linie fällt
ins Sichtbare, was für Li$^{2+}$ typisch ist.
\end{kontrolle}

\begin{merksatz}
Wasserstoffähnlich = Formel mit $Z^2$-Faktor; aus $\nu$ folgen
$\lambda = c/\nu$ und $\tilde\nu = 1/\lambda$.
\end{merksatz}
